% ==================== 脚注编号、字体、分隔线与间距 ====================

%%%%%%%% 脚注 %%%%%%%%

% 按实际输出页重置；需要复编后才能稳定处理自动分页附近的脚注。
\RequirePackage{perpage}
% 每页重置取代 ctexbook 的按章重置，避免章节开始时产生一次虚拟脚注计数。
\@removefromreset{footnote}{chapter}
\MakePerPage{footnote}

% 脚注字体
% 圆圈显式使用符号字体，数字仍为 Times New Roman，避免 OMS/Times 字形回退。
\newcommand{\xmu@circlednumber}[1]{%
    \begingroup\fontencoding{OMS}\fontfamily{cmsy}\selectfont
        \textcircled{\normalfont\rmfamily\zihao{7}#1}%
    \endgroup
}
\renewcommand{\thefootnote}{\xmu@circlednumber{\arabic{footnote}}}
\renewcommand{\thempfootnote}{\xmu@circlednumber{\arabic{mpfootnote}}}
\renewcommand{\footnotesize}{\linespread{1.5}\xmu@songti\zihao{-5}}
% 脚注线
\renewcommand{\footnoterule}{\vspace{1em}\hrule width 0.4\columnwidth height 0.4pt\vspace{0.3em}}
% 脚注缩进、行距
\renewcommand{\footnotesep}{1.2em}
\deffootnote[3ex]{3ex}{1em}{\zihao{7}\raisebox{0.4ex}{\rm\zihao{-5}\thefootnotemark}}
\AddToHook{cmd/@makefntext/before}{\raggedright}
